\section{À quel problème répond votre projet ?}

% ----------------------------------------------------------------
% 📌 À quoi sert cette section ?
% ----------------------------------------------------------------
% Cette partie permet de contextualiser l'objectif du projet en expliquant :
% - Le problème initial qui a motivé sa création
% - Les limites ou lacunes des solutions existantes
% - Les enjeux techniques, humains ou pédagogiques concernés
%
% C’est une section essentielle qui justifie l’intérêt du projet.

% ----------------------------------------------------------------
% 🧭 Comment la remplir ?
% ----------------------------------------------------------------
% ➤ Identifiez clairement un problème ou une frustration vécue par les utilisateurs ou les acteurs concernés.
% ➤ Expliquez en quoi ce problème est réel, fréquent ou bloquant.
% ➤ Si possible, illustrez avec des cas d’usage, des témoignages, ou des situations observées.
% ➤ Terminez par une phrase de transition vers la solution proposée dans le projet.

% ----------------------------------------------------------------
% 🧾 Exemple rédigé :
% ----------------------------------------------------------------
% Aujourd’hui, l’apprentissage du format SVG reste très théorique dans les cursus de développement front-end.
% Les ressources disponibles sont nombreuses mais peu interactives et rarement adaptées aux débutants.
% Il n’existe pas de plateforme ludique francophone permettant d’expérimenter directement avec le code SVG
% sans configuration complexe.
%
% Ce manque de support pratique freine l’autonomisation des étudiants et limite leur motivation.
% Le projet vise à combler cette lacune avec un outil visuel accessible en ligne.

% ----------------------------------------------------------------
% 📊 Exemple de schéma à insérer (problème & causes)
% ----------------------------------------------------------------

% \begin{figure}[H]
%     \centering
%     \begin{tikzpicture}[node distance=1.5cm]
%         \node (pb) [draw, rectangle, fill=red!10] {Problème : apprentissage SVG non pratique};
%         \node (c1) [below left=of pb, draw, rectangle] {Absence d'outils interactifs};
%         \node (c2) [below right=of pb, draw, rectangle] {Ressources trop complexes};
%         \draw[->] (c1) -- (pb);
%         \draw[->] (c2) -- (pb);
%     \end{tikzpicture}
%     \caption{Causes du problème identifié}
%     \label{fig:problematique}
% \end{figure}

% ----------------------------------------------------------------
% ✍️ À vous de jouer : complétez ci-dessous
% ----------------------------------------------------------------

% (Exemple de départ)
Le problème identifié dans le cadre de ce projet est que...

